\documentclass[a4paper,12pt,twoside]{memoir}

% Castellano
\usepackage[spanish,es-tabla]{babel}
\selectlanguage{spanish}
\usepackage[utf8]{inputenc}
\usepackage[T1]{fontenc}
\usepackage{lmodern} % Scalable font
\usepackage{microtype}
\usepackage{placeins}

\RequirePackage{booktabs}
\RequirePackage[table]{xcolor}
\RequirePackage{xtab}
\RequirePackage{multirow}

% Links
\usepackage[colorlinks]{hyperref}
\hypersetup{
	allcolors = {red}
}

% Ecuaciones
\usepackage{amsmath}

% Rutas de fichero / paquete
\newcommand{\ruta}[1]{{\sffamily #1}}

% Párrafos
\nonzeroparskip


% Imagenes
\usepackage{graphicx}
\newcommand{\imagen}[2]{
	\begin{figure}[!h]
		\centering
		\includegraphics[width=0.9\textwidth]{#1}
		\caption{#2}\label{fig:#1}
	\end{figure}
	\FloatBarrier
}

\newcommand{\imagenflotante}[2]{
	\begin{figure}%[!h]
		\centering
		\includegraphics[width=0.9\textwidth]{#1}
		\caption{#2}\label{fig:#1}
	\end{figure}
}



% El comando \figura nos permite insertar figuras comodamente, y utilizando
% siempre el mismo formato. Los parametros son:
% 1 -> Porcentaje del ancho de página que ocupará la figura (de 0 a 1)
% 2 --> Fichero de la imagen
% 3 --> Texto a pie de imagen
% 4 --> Etiqueta (label) para referencias
% 5 --> Opciones que queramos pasarle al \includegraphics
% 6 --> Opciones de posicionamiento a pasarle a \begin{figure}
\newcommand{\figuraConPosicion}[6]{%
  \setlength{\anchoFloat}{#1\textwidth}%
  \addtolength{\anchoFloat}{-4\fboxsep}%
  \setlength{\anchoFigura}{\anchoFloat}%
  \begin{figure}[#6]
    \begin{center}%
      \Ovalbox{%
        \begin{minipage}{\anchoFloat}%
          \begin{center}%
            \includegraphics[width=\anchoFigura,#5]{#2}%
            \caption{#3}%
            \label{#4}%
          \end{center}%
        \end{minipage}
      }%
    \end{center}%
  \end{figure}%
}

%
% Comando para incluir imágenes en formato apaisado (sin marco).
\newcommand{\figuraApaisadaSinMarco}[5]{%
  \begin{figure}%
    \begin{center}%
    \includegraphics[angle=90,height=#1\textheight,#5]{#2}%
    \caption{#3}%
    \label{#4}%
    \end{center}%
  \end{figure}%
}
% Para las tablas
\newcommand{\otoprule}{\midrule [\heavyrulewidth]}
%
% Nuevo comando para tablas pequeñas (menos de una página).
\newcommand{\tablaSmall}[5]{%
 \begin{table}
  \begin{center}
   \rowcolors {2}{gray!35}{}
   \begin{tabular}{#2}
    \toprule
    #4
    \otoprule
    #5
    \bottomrule
   \end{tabular}
   \caption{#1}
   \label{tabla:#3}
  \end{center}
 \end{table}
}

%
% Nuevo comando para tablas pequeñas (menos de una página).
\newcommand{\tablaSmallSinColores}[5]{%
 \begin{table}[H]
  \begin{center}
   \begin{tabular}{#2}
    \toprule
    #4
    \otoprule
    #5
    \bottomrule
   \end{tabular}
   \caption{#1}
   \label{tabla:#3}
  \end{center}
 \end{table}
}

\newcommand{\tablaApaisadaSmall}[5]{%
\begin{landscape}
  \begin{table}
   \begin{center}
    \rowcolors {2}{gray!35}{}
    \begin{tabular}{#2}
     \toprule
     #4
     \otoprule
     #5
     \bottomrule
    \end{tabular}
    \caption{#1}
    \label{tabla:#3}
   \end{center}
  \end{table}
\end{landscape}
}

%
% Nuevo comando para tablas grandes con cabecera y filas alternas coloreadas en gris.
\newcommand{\tabla}[6]{%
  \begin{center}
    \tablefirsthead{
      \toprule
      #5
      \otoprule
    }
    \tablehead{
      \multicolumn{#3}{l}{\small\sl continúa desde la página anterior}\\
      \toprule
      #5
      \otoprule
    }
    \tabletail{
      \hline
      \multicolumn{#3}{r}{\small\sl continúa en la página siguiente}\\
    }
    \tablelasttail{
      \hline
    }
    \bottomcaption{#1}
    \rowcolors {2}{gray!35}{}
    \begin{xtabular}{#2}
      #6
      \bottomrule
    \end{xtabular}
    \label{tabla:#4}
  \end{center}
}

%
% Nuevo comando para tablas grandes con cabecera.
\newcommand{\tablaSinColores}[6]{%
  \begin{center}
    \tablefirsthead{
      \toprule
      #5
      \otoprule
    }
    \tablehead{
      \multicolumn{#3}{l}{\small\sl continúa desde la página anterior}\\
      \toprule
      #5
      \otoprule
    }
    \tabletail{
      \hline
      \multicolumn{#3}{r}{\small\sl continúa en la página siguiente}\\
    }
    \tablelasttail{
      \hline
    }
    \bottomcaption{#1}
    \begin{xtabular}{#2}
      #6
      \bottomrule
    \end{xtabular}
    \label{tabla:#4}
  \end{center}
}

%
% Nuevo comando para tablas grandes sin cabecera.
\newcommand{\tablaSinCabecera}[5]{%
  \begin{center}
    \tablefirsthead{
      \toprule
    }
    \tablehead{
      \multicolumn{#3}{l}{\small\sl continúa desde la página anterior}\\
      \hline
    }
    \tabletail{
      \hline
      \multicolumn{#3}{r}{\small\sl continúa en la página siguiente}\\
    }
    \tablelasttail{
      \hline
    }
    \bottomcaption{#1}
  \begin{xtabular}{#2}
    #5
   \bottomrule
  \end{xtabular}
  \label{tabla:#4}
  \end{center}
}



\definecolor{cgoLight}{HTML}{EEEEEE}
\definecolor{cgoExtralight}{HTML}{FFFFFF}

%
% Nuevo comando para tablas grandes sin cabecera.
\newcommand{\tablaSinCabeceraConBandas}[5]{%
  \begin{center}
    \tablefirsthead{
      \toprule
    }
    \tablehead{
      \multicolumn{#3}{l}{\small\sl continúa desde la página anterior}\\
      \hline
    }
    \tabletail{
      \hline
      \multicolumn{#3}{r}{\small\sl continúa en la página siguiente}\\
    }
    \tablelasttail{
      \hline
    }
    \bottomcaption{#1}
    \rowcolors[]{1}{cgoExtralight}{cgoLight}

  \begin{xtabular}{#2}
    #5
   \bottomrule
  \end{xtabular}
  \label{tabla:#4}
  \end{center}
}


















\graphicspath{ {./img/} }

% Capítulos
\chapterstyle{bianchi}
\newcommand{\capitulo}[2]{
	\setcounter{chapter}{#1}
	\setcounter{section}{0}
	\chapter*{#2}
	\addcontentsline{toc}{chapter}{#2}
	\markboth{#2}{#2}
}

% Apéndices
\renewcommand{\appendixname}{Apéndice}
\renewcommand*\cftappendixname{\appendixname}

\newcommand{\apendice}[1]{
	%\renewcommand{\thechapter}{A}
	\chapter{#1}
}

\renewcommand*\cftappendixname{\appendixname\ }

% Formato de portada
\makeatletter
\usepackage{xcolor}
\newcommand{\tutor}[1]{\def\@tutor{#1}}
\newcommand{\course}[1]{\def\@course{#1}}
\definecolor{cpardoBox}{HTML}{E6E6FF}
\def\maketitle{
  \null
  \thispagestyle{empty}
  % Cabecera ----------------
\noindent\includegraphics[width=\textwidth]{cabecera}\vspace{1cm}%
  \vfill
  % Título proyecto y escudo informática ----------------
  \colorbox{cpardoBox}{%
    \begin{minipage}{.8\textwidth}
      \vspace{.5cm}\Large
      \begin{center}
      \textbf{TFG del Grado en Ingeniería Informática}\vspace{.6cm}\\
      \textbf{\LARGE\@title{}}
      \end{center}
      \vspace{.2cm}
    \end{minipage}

  }%
  \hfill\begin{minipage}{.20\textwidth}
    \includegraphics[width=\textwidth]{escudoInfor}
  \end{minipage}
  \vfill
  % Datos de alumno, curso y tutores ------------------
  \begin{center}%
  {%
    \noindent\LARGE
    Presentado por \@author{}\\ 
    en Universidad de Burgos --- \@date{}\\
    Tutores: \@tutor{}\\
  }%
  \end{center}%
  \null
  \cleardoublepage
  }
\makeatother

\newcommand{\nombre}{Adrián Marcos Batlle} %%% cambio de comando

% Datos de portada
\title{{\Huge CRET}\\[0.5cm]CanBus Reverse Engineering Toolkit}
\author{\nombre}
\tutor{Álvar Arnaiz González y César Represa Pérez}
\date{\today}

\begin{document}

\maketitle


\newpage\null\thispagestyle{empty}\newpage


%%%%%%%%%%%%%%%%%%%%%%%%%%%%%%%%%%%%%%%%%%%%%%%%%%%%%%%%%%%%%%%%%%%%%%%%%%%%%%%%%%%%%%%%
\thispagestyle{empty}


\noindent\includegraphics[width=\textwidth]{cabecera}\vspace{1cm}

\noindent D. Álvar Arnaiz González profesor del departamento de Ingeniería Civil, área de Lenguajes y Sistemas Informáticos, y D. César Represa Pérez, profesor del departamento de Ingeniería Electromecánica.

\noindent Exponen:

\noindent Que el alumno D. \nombre, con DNI -, ha realizado el Trabajo final de Grado en Ingeniería Informática titulado CRET - CanBus Reverse Engineering Toolkit. 

\noindent Y que dicho trabajo ha sido realizado por el alumno bajo la dirección del que suscribe, en virtud de lo cual se autoriza su presentación y defensa.

\begin{center} %\large
En Burgos, {\large \today}
\end{center}

\vfill\vfill\vfill

% Author and supervisor
\begin{minipage}{0.45\textwidth}
\begin{flushleft} %\large
Vº. Bº. del Tutor:\\[2cm]
D. Álvar Arnaiz González
\end{flushleft}
\end{minipage}
\hfill
\begin{minipage}{0.45\textwidth}
\begin{flushleft} %\large
Vº. Bº. del co-tutor:\\[2cm]
D. César Represa Pérez
\end{flushleft}
\end{minipage}
\hfill

\vfill

% para casos con solo un tutor comentar lo anterior
% y descomentar lo siguiente
%Vº. Bº. del Tutor:\\[2cm]
%D. nombre tutor


\newpage\null\thispagestyle{empty}\newpage


\frontmatter

% Abstract en castellano
\renewcommand*\abstractname{Resumen}
\begin{abstract}

El bus CAN (\emph{CAN Bus - Controlled Area Network}) es un protocolo de comunicación utilizado en los desarrollos de multitud de sectores críticos como la industria, la automoción y la aviación entre otros, para la comunicación entre los componentes internos que forman la infraestructura desarrollada.

Los datos que son transmitidos por dicho bus son propiedad de cada uno de los fabricante, a pesar de ser un protocolo con los estándares públicos.

El desarrollo de este proyecto viene motivado a realizar un análisis de la información que circula por estos buses, así como su clasificación y monitorización en tiempo real.

Para dicho objetivo, se ha desarrollado tanto una parte de \emph{software} (para el análisis, clasificación y monitorización de los datos), como una parte de \emph{hardware} necesaria para conectarse a dicho bus.

Con ello, se persigue el objetivo de ahorrar tiempo y costes, por ejemplo, en caso de que necesitemos obtener una señal de un sensor conectado a dicha red. En lugar de instalar otro sensor nuevo, sería posible obtener el valor del ya instalado una vez se haya identificado.

\end{abstract}

\renewcommand*\abstractname{Descriptores}
\begin{abstract}
Bus CAN, Automoción, Industria, Aviónica, Ingeniería inversa, Análisis de protocolos.
\end{abstract}

\clearpage

% Abstract en inglés
\renewcommand*\abstractname{Abstract}
\begin{abstract}
CAN Bus (\emph{Controlled Area Network}) is a communication protocol used in the developement of a large number of critical sectors such as industry, automotive and aviation among others, for communication between internal electronic components.

Data transmited throught the CAN bus is property of each developer, despite being a protocol in which the standards are public.

The developement of this proyect is motivated to perform an analysis of the data that flows througth these buses, as well as their classification and real time monitoring.

For this purpose, both custom software and hardware have been developed to achieve the requirements.

This means we can save time and costs, for example, getting the signal of a sensor connected to the bus, instead of installing a new one.
\end{abstract}

\renewcommand*\abstractname{Keywords}
\begin{abstract}
Can Bus, Automotive, Industry, Avionics, Reversing, Protocol analysis.
\end{abstract}

\clearpage

% Indices
\tableofcontents

\clearpage

\listoffigures

\clearpage

\listoftables
\clearpage

\mainmatter
\capitulo{1}{Introducción}

Descripción del contenido del trabajo y del estrucutra de la memoria y del resto de materiales entregados.

\capitulo{2}{Objetivos del proyecto}

Este apartado explica de forma precisa y concisa cuales son los objetivos que se persiguen con la realización del proyecto. Se puede distinguir entre los objetivos marcados por los requisitos del software a construir y los objetivos de carácter técnico que plantea a la hora de llevar a la práctica el proyecto.



\section{Objetivos generales}\label{objetivos-generales}

\begin{itemize}
\tightlist
\item
  Desarrollar una aplicación para el análisis y monitorización de las señales que fluyen por el bus CAN.
  
\end{itemize}

\section{Objetivos técnicos}\label{objetivos-tecnicos}

\begin{itemize}
\tightlist
\item
  Desarrollar un hardware propio desde 0, siguiendo la metodología para el desarrollo del mismo.
\item
  Desarrollar una aplicación en Java y JavaFX 
\item
  Aplicar la arquitectura MVP (\emph{Model-View-Presenter}) en el
  desarrollo de la aplicación.
\item
  Utilizar Zenhub (basado en el método Kanban) para realizar un seguimiento y gestión de proyectos.
\item
  Utilizar Git (en la plataforma GitHub) para realizar un control de versiones de software.
\end{itemize}

\section{Objetivos personales}\label{objetivos-personales}

\begin{itemize}
\tightlist
\item
  Realizar una aportación a la modernización de la apicultura.

\end{itemize}
\capitulo{3}{Conceptos teóricos}

Las partes del proyecto con mayor desconocimiento y complejidad están enfocadas principalmente en el funcionamiento del bus CAN y en el desarrollo del \emph{hardware}, el cual requiere de unos conocimientos básicos y unas metodologías específicas las cuales serán detalladas a continuación:

\section{Bus CAN}\label{bus_can}

El bus CAN (\emph{Controller Area Network})\cite{wiki:can} es un protocolo desarrollado para la comunicación entre los distintos micro-controladores y dispositivos que son necesarios para el funcionamiento de una máquina (un vehículo, por ejemplo). Este protocolo no necesita de un host principal, sino que sigue una topología de tipo "bus".

\imagen{canNetwork}{Ejemplo de dispositivos conectados a un bus CAN.}

Dicho protocolo está basado en el uso de mensajes para el intercambio de información entre los distintos dispositivos que lo componen.

Es utilizado en multitud de escenarios como la aviación, la navegación, la automatización industrial, instrumentos médicos, maquinaria pesada y ascensores, entre otros.

Una de las características que nos interesa conocer, es que al ser una topología de tipo BUS, todos los nodos tienen acceso a la información que es transmitida a través del bus, con lo que es suficiente con conectarnos a dicho bus y capturar los datos que viajan a través de él para analizarlos posteriormente.

\imagen{canbusTopology}{Esquema de conexión de los nodos en el bus CAN.}

El principal elemento que compone una red de tipo CAN son los nodos, además de los cables de conexión entre los mismos:

\begin{itemize}
\item
\textbf{Nodo:} Cada uno de los dispositivos físicos que están conectados a la red CAN. Al menos es necesario que existan dos nodos conectados a la red para que se produzca una comunicación.
\item
\textbf{Cable:} Se trata del cable a través del cual se envían y reciben los datos. Todos los nodos deben de estar conectados a estos dos cables. Estos son identificados por CAN-H (\emph{CAN-High}) y CAN-L (\emph{CAN-Low}).
\end{itemize}

Esto es debido a que las señales que circulan por los mismos, son señales diferenciales, de las cuales el valor efectivo interpretado por el nodo es extraído tras aplicar la diferencia entre ambas como se puede ver en la figura.

\imagen{canSignals}{Señales eléctricas del bus CAN vistas por un osciloscopio.\cite{marco:can}}

Cada uno de los nodos es capaz de enviar y recibir mensajes. La prioridad de los mensajes de un nodo viene dada por la ID(identificador) del \emph{Frame} enviado al bus.

\imagen{canNode}{Anatomía de un nodo CAN.}

A continuación se describen los elementos de una trama CAN, así como la descripción de la misma.

\textbf{\emph{Frame} (Trama):} Cada uno de los mensajes enviados a través de la red CAN.

La estructura y el funcionamiento interno es más complejo, pero no nos centraremos en ellos ya que no tiene relevancia en este caso. Únicamente nos centraremos en los tres elementos identificados en la figura.

Cada uno de los \emph{Frames} contiene tres campos en los que nos centraremos: ID, Length y Data.

\begin{itemize}
\item
\textbf{ID (Identificador):} Se trata del identificador del \emph{Frame}. A través de él, se establece la prioridad que tiene ese \emph{Frame} durante el acceso al bus de datos.
\item
\textbf{Length (Longitud):} En este campo, se define la longitud del campo de datos del \emph{Frame}. En las versiones estándar, esta longitud puede ir desde 0 hasta 8 \emph{bytes} de datos.
\item
\textbf{Data (Datos):} Campo en el que se transportan los datos que son enviados por el nodo correspondiente. La longitud de este campo puede ir desde 0 \emph{bytes} hasta un total de 8 \emph{bytes}.
\end{itemize}

\imagen{candump}{Identificación de los campos que componen un \emph{Frame}.}


\textbf{Bitrate:} Se trata de la velocidad a la que los datos son transmitidos a través del bus. Todos los nodos deben de transmitir a la misma velocidad para entenderse entre ellos. Esta, puede cambiar dependiendo de los sistemas y de su desarrollo.
Los valores más comunes son:

\begin{itemize}
\tightlist
\item
10\,000 bit/s
\item
20\,000 bit/s
\item
50\,000 bit/s
\item
100\,000 bit/s
\item
125\,000 bit/s
\item
250\,000 bit/s
\item
500\,000 bit/s
\item
800\,000 bit/s
\item
1\,000\,000 bit/s
\end{itemize}


Además, el estándar dispone de varias medidas para la detección de errores y seguridad las cuales no son relevantes para la exposición de este proyecto.

El desarrollo del \emph{hardware} con dos interfaces viene motivado a la posibilidad de que exista más de un bus en una misma máquina. De esta manera, sería posible el análisis de dos buses de forma simultánea.

Siguiendo el modelo OSI\cite{osi:can}, el estándar bus CAN especifica únicamente las dos primeras capas:

\begin{itemize}

\item
\textbf{Capa física:} Esta capa define cómo son transmitidas las señales a nivel eléctrico, es decir, los niveles de las señales en su representación como bits, y el medio de transmisión que va a ser utilizado. Además, se encargaría de gestionar la sincronización de los mensajes, así como su \emph{bit encoding}.

\item
\textbf{Capa de enlace de datos:} Esta capa se encarga de la parte lógica del protocolo. Entre sus funciones se encuentra el filtrado de los \emph{Frames}, detección de sobrecarga del bus y la detección de errores.

\end{itemize}

Existen 3 modos por los cuales los nodos pueden conectarse al bus CAN:

\begin{itemize}
\item
\textbf{Active:} Cuando un nodo es conectado en modo \emph{active}, significa que si envía un mensaje a través del bus. Al no ser que el \emph{transciver} esté conectado en modo \emph{Single Shot}, el mensaje será enviado continuamente por el bus hasta que se reciba un ACK \emph{(Acknonolowment)}. De manera que este modo si que tiene influencias sobre el bus.

\item
\textbf{Listen-Only:} Este modo no interfiere en el funcionamiento del bus CAN. Cuando un nodo es conectado en modo \emph{listenonly}, dicho nodo no tiene influencia sobre el resto. Si por ejemplo llega un \emph{Frame} erróneo, no genera \emph{flags} de error. Se podría decir que es un modo de escucha del bus, sin posibilidad de interacción con el mismo.

\item
\textbf{LoopBack:} Se podría considerar una especie de modo \emph{debug}. Cuando un nodo se conecta en modo \emph{LoopBack}, significa que este nodo es capaz comunicarse consigo mismo, sin enviar ni recibir mensajes del bus. Es un modo utilizado solo para la realización de pruebas, no para aplicaciones finales.

\end{itemize}


\section{Desarrollo del hardware}\label{desarrollo_del_hardware}

A continuación se describen distintos conceptos básicos sobre el \emph{hardware} desarrollado y algunos de sus componentes:

\subsection{Esquema}\label{esquema}

Representación esquemática de los componentes de un circuito, así como sus entradas y salidas, y conexión entre ellos.

A continuación se muestra un ejemplo visual de varios componentes en un esquema para el diseño de una PCB\emph{(Printed Circuit Board)}:

\imagen{esquema}{Ejemplo de varios componentes a nivel esquemático.}

\subsection{PCB - Printed Circuit Board}\label{pcb_printed_circuit_board}

Hace referencia a la placa sobre la que van soldados todos los componentes eléctricos necesarios para hacer funcionar el circuito. Suele tener 3 componentes principales:

\begin{itemize}
\item
Pistas
\item
\emph{Pads}
\item
Vías
\end{itemize}
\newpage
\subsection{Pistas}\label{pistas}

Cada una de las conexiones entre los componentes. Hacen el trabajo de un cable, solo que integrado en la placa.

\imagen{pistas}{Esquema de pistas de una PCB.}

\subsection{Pad}\label{pad}

Cada uno de los espacios en los cuales se sueldan los distintos  elementos que componen el circuito.

\imagen{pad}{Ilustración con 4 \emph{pads} en una PCB.}

\subsection{Vía}\label{via}

Al tratarse de una PCB de doble capa (es decir, que tiene pistas o elementos en ambos lados), a veces es necesario realizar un cruce entre vías, o conectar un componente que está en el otro lado de la placa. Para ello se utilizan las vías, a través de ellas se consiguen conectar ambas capas de la PCB.

\imagen{via}{Visión de una vía a través de un corte de una PCB.}

\subsection{Footprint}\label{footprint}

Hace referencia al espacio necesario para colocar un elemento del \emph{hardware}. Existen estándares que definen los tamaños y formas de dichos espacios.

\imagen{footprint}{Ejemplo de una huella de un IC \emph{(Integrated Circuit)}.}

\subsection{Gerber}\label{gerber}
Es un formato de fichero el cual contiene la información necesaria para realizar la fabricación de una PCB.

Estos ficheros son aceptados por todos los fabricantes de placas de circuito impreso. Además es soportado por la mayoría de programas utilizados para el diseño de las mismas.

\section{Desarrollo del software}\label{desarrollo_del_software}


Para el desarrollo del \emph{software} se ha utilizado el lenguaje de programación Java. Además, para la mejora de la interacción con el usuario también se ha utilizado el módulo JavaFX y la librería \emph{medusa}, un complemento para JavaFX.

Existen distintos elementos utilizados durante el desarrollo de la aplicación los cuales es interesante destacar:

\begin{itemize}
\item
\textbf{GridPane}: Se trata de una estructura implementada en JavaFX la cual nos permite una distribución de los elementos que forman la interfaz gráfica en forma de cuadrícula.

En cada una de las celdas de esta cuadrícula se incluyen distintos elementos que se mencionarán a continuación.
\item
\textbf{LineChart}: Se trata de un módulo perteneciente a JavaFX el cual nos permite generar una gráfica. Es utilizado para mostrar los datos al usuario y que estos puedan ser identificados.

\emph{LineChart} recoge los datos de una estructura de tipo \emph{ObservableList}, en la cual se encuentran todos los datos a representar.

Este tipo de datos no permite el acceso concurrente de varios \emph{threads} al mismo tiempo.

\imagen{lineChart}{Ejemplo del tipo \emph{LineChart} en JavaFX.}

\item
\textbf{ConcurrentLinkedQueue}: A través del uso de esta estructura ha sido posible recoger los datos del bus CAN con un \emph{thread} a la vez que el \emph{thread} de la interfaz gráfica se encargaba de insertarlos en los objetos de tipo \emph{ObservableList}.
\item
\textbf{FXML Document}: Se trata de documentos con un formato que extiende de XML. En estos documentos se encuentra la estructura de las interfaces gráficas y la distribución de los elementos en la misma. Este tipo de documentos pueden ser escritos a mano o generados con la ayuda de la herramienta \emph{SceneBuilder}.

A este documento se le indica el \emph{Controller} que va a manejarlo, de forma que se crea la conexión \emph{View-Controller}. Esto se hace a través del campo \emph{fx:controller} del documento FXML.
\item
\textbf{Preloader}: Debido a la cantidad de tiempo requerido para el arranque de la aplicación (comprobación de la base de datos, preparación de la interfaz gráfica), se decidió introducir un \emph{preloader}. Este se encarga de mostrarnos una pantalla de carga mientras la aplicación realiza las tareas necesarias. En el momento en el que la pantalla principal está lista para mostrarse, esta manda una señal al \emph{preloader} para que se cierre.

\imagen{preloaderArchitecture}{Etapas del arranque de una aplicación JavaFX.\cite{preloader:javafx}}
\item
\textbf{ToolTips}: Se trata de una descripción emergente de un elemento. Esta ayuda al usuario a entender la funcionalidad de un elemento simplemente arrastrando el cursor sobre el elemento.

\imagen{tooltip}{Ejemplo de ToolTip en JavaFX.}

\end{itemize}



\capitulo{4}{Técnicas y herramientas}

Esta parte de la memoria tiene como objetivo presentar las técnicas metodológicas y las herramientas de desarrollo que se han utilizado para llevar a cabo el proyecto. Si se han estudiado diferentes alternativas de metodologías, herramientas, bibliotecas se puede hacer un resumen de los aspectos más destacados de cada alternativa, incluyendo comparativas entre las distintas opciones y una justificación de las elecciones realizadas. 
No se pretende que este apartado se convierta en un capítulo de un libro dedicado a cada una de las alternativas, sino comentar los aspectos más destacados de cada opción, con un repaso somero a los fundamentos esenciales y referencias bibliográficas para que el lector pueda ampliar su conocimiento sobre el tema.



\capitulo{5}{Aspectos relevantes del desarrollo del proyecto}

A continuación se desarrollan los aspectos más importantes que influyeron en el desarrollo del proyecto.

\section{Inicio del proyecto}\label{inicio_del_proyecto}

La idea de desarrollar esta aplicación surgió unos meses atrás, tras trabajar en una empresa que de dedicaba a la reparación de maquinaria.

Tras investigar y leer algunos manuales de reparación de esas máquinas, era de destacar que existía una característica en común, que era el uso del bus CAN en muchas de ellas, por no decir todas. A partir de ahí surgió el interés y las ganas de indagar sobre su funcionamiento.

Una vez fue posible la lectura de dichos datos, quedaba clara la dificultad para la identificación de los mismos, ya que todos ellos estaban en hexadecimal, lo cual dificultaba su identificación.

Tras el desarrollo de un pequeños \emph{script} en python los cuales eran capaces de realizar una gráfica con los datos previamente adquiridos del bus CAN y detectar la mayor facilidad a la hora de identificar esos datos, surgió la idea de realizar este proyecto.

Existía la necesidad de que, al contrario que en el caso anterior, este proyecto permitiera realizar ese proceso en tiempo real, no con una muestra de datos tomada anteriormente.

Otra de las principales motivaciones fue el hecho de que si por ejemplo, existiera la necesidad de recoger los datos de un sensor (de velocidad de una rueda) para su monitorización y este dato es transmitido a través del bus CAN, no es necesaria la instalación de un nuevo sensor, abaratando costes y facilitando realizar la tarea necesaria.

\section{Metodologías}\label{metodologias}

Se decidió utilizar la metodología \emph{Scrum}, una metodología ágil, de forma que no era necesario definir el proyecto desde el principio (no se conocía el alcance del mismo, ni si iba a ser posible realizar ciertas propuestas planteadas).

La duración aproximada de los \emph{sprints} era de una semana, a la vez que al finalizar dicho \emph{sprint}, se planificaba el de la semana siguiente.

A través del uso de la herramienta \emph{ZenHub}, se generaban \emph{issues}, es decir, tareas a realizar en ese \emph{sprint}, y se cambiaban de estado según se iban finalizando.

\section{Formación}\label{formacion}

El proyecto requería conceptos tanto de funcionamiento como de la implementación del bus CAN, los cuales desconocía en un principio. A través de los siguientes documentos, se adquirieron conocimientos sobre la materia:


LINKS CAN BUS


\section{Desarrollo del \emph{hardware}}\label{desarrollo_del_hardware}

Se decidió comenzar por esta fase por los siguientes motivos:

\begin{itemize}
\item
No era seguro que el \emph{hardware} fuera a funcionar con lo que no se podía crear una dependencia por parte del software en caso de que esa parte no se realizase satisfactoriamente.

\item
La producción de la PCB requiere de unas semanas desde que envías el pedido hasta que el material llega.
\item
Era la parte más desconocida y en la cual habría que invertir más tiempo para aplicar la metodología sugerida a la hora del desarrollo de este tipo de proyectos.
\end{itemize}

Se podrían definir 4 etapas las cuales se han llevado a cabo para el desarrollo:

\subsection{Diseño del esquema eléctrico}\label{diseño_del_esquema_electrico}

El primer paso para el desarrollo de una placa de circuito impreso, es diseñar el esquema en el que se incluyen los componentes, y la conexión que se va a realizar entre ellos. En este caso, el objetivo era duplicar un esquema existente, introduciéndole un \emph{Hub USB}, todo ello en la misma placa.

En este esquema se especifican los componentes utilizados (resistencias, condensadores, circuitos integrados, etc). Además se define que entradas y salidas de esos componentes van conectados al resto de componentes.

\imagen{scheme}{Vista general del esquema.}

\subsection{Generación del \emph{netlist}}\

Una vez realizado el esquema eléctrico, necesitamos generar el fichero \emph{Netlist}. En este fichero se almacenan los siguientes datos:

Por cada componente:

\begin{itemize}
\item
Se describen las conexiones de cada uno de los elementos del circuito.
\item
El tipo de elemento que es cada componente.
\item
La huella (\emph{footprint}) que va a utilizar, es decir, en función de cada encapsulado (forma física del \emph{chip}), será necesario que la zona de soldadura de la PCB sea adaptada al mismo, tanto en tamaño como en forma.
\end{itemize}


\subsection{Posicionamiento y enrutado de los componentes y las pistas en la PCB}\label{posicionamiento_y_enrutado_de_los_componentes_y_las_pistas_en_la_pcb}

Una vez se ha generado el fichero \emph{Netlist}, este es cargado en el editor de PCBs. Para que nos hagamos una idea, ese es el resultado:


En este momento, es cuando tenemos que decidir:

\begin{itemize}
\item
El número de capas que va a tener nuestra PCB: En este caso se utilizaron 2 capas al tratarse de un proyecto sencillo, pero en el que necesitaban hacerse cruces de pistas.
\item
La posición de los componentes y el enrutado de las pistas.
\end{itemize}

FOTO DEL ESQUEMA TERMINADO.


\imagen{fotoPCB}{PCB del proyecto sin los componentes}


Tras limpiar la PCB, posicionar y soldar los componentes, obtenemos el siguiente resultado:


\imagen{finishedPCB}{PCB del proyecto con los componentes soldados}

\section{Desarrollo del \emph{software}}

Lo primero que había que decidir era el lenguaje en el que se iba a desarrollar la aplicación. Se barajó la posibilidad de realizarlo en Java o Python. La decisión de Java se debió a la disposición de una librería compatible con el \emph{hardware} y a la librería JavaFX la cual nos permite realizar una aplicación muy visual y amigable de forma sencilla y eficiente.

Una vez elegido el lenguaje de desarrollo, se comenzó con la modelación de la interfaz gráfica, es decir, de la estructura de la aplicación a nivel visual. Una vez estaba claro, se comenzó con el desarrollo del código.

Para la base de datos, se decidió utilizar SQLite, ya que no requiere de ningún motor de base de datos, sino simplemente la utilización de una librería interna en la aplicación. Además, la información almacenada iba a ser muy pequeña y era más que suficiente.

Principalmente surgieron dos grandes problemas durante el desarrollo de la parte del software, ambos relacionados con la parte de la interfaz gráfica.

\subsection{Generación de gráficas de forma dinámica}\label{generacion_de_graficas_de_forma_dinamica}

Era necesario que las gráficas se fueran creando de manera dinámica cuando existiera un nuevo dato a representar, no tener una estructura estática creada previamente. Para ello, cada vez que un nuevo \emph{Frame} del bus CAN fuera detectado, había que proceder a la creación de dichos elementos, así como de un \emph{listener} en ese objeto, el cual consultase la estructura en la que son almacenados los datos, para su posterior representación gráfica.

Para ello se decidió utilizar una estructura gráfica de tipo \emph{GridPane}, dividiendo la pantalla en el eje horizontal en 3 celdas. Cada una de estas celdas se iría rellenando en función de las necesidades e iría incrementandose de arriba a abajo si fuera necesario. Para conseguir ese efecto, el elemento padre del \emph{GridPane} tenía que ser un \emph{ScrollPane} es decir, lo que conocemos como un \emph{Scroll}.

Todo esto tenía que ir de la mano con una serie de estructuras, también generadas dinámicamente, para realizar la lectura y escritura de los datos que tienen que ser representados.

De esa manera, se decidió utilizar estructuras de tipo \emph{HashMap}, las cuales identificaban el dato por el par \emph{ID} y \emph{número de byte representado}.

\subsection{Problemas de concurrencia}\label{problemas_de_concurrencia}

Los elementos que representan gráficamente los datos (LineChart), realizan una lectura de un tipo de datos llamado \emph{ObservableList}. Durante la ejecución de la aplicación existen 2 \emph{threads}, uno para la ejecución de la parte gráfica, y el otro para la lectura de los datos del bus CAN.

El problema surge cuando se intenta realizar una escritura en el objeto de tipo \emph{ObservableList} desde el \emph{thread} que lee el bus CAN al mismo tiempo que el \emph{thread} de la interfaz gráfica realiza una lectura de los datos del mismo objeto.

Para la resolución de este problema se decidió utilizar un modelo de tipo Productor/Consumidor. Para ello se definió un objeto de tipo \emph{ConcurrentLinkedQueue} en el cual el \emph{thread} encargado de la lectura del bus CAN, iba introduciendo la información que le llegaba al tiempo que el \emph{thread} de la interfaz gráfica iba cargando esos datos en el objeto de tipo \emph{ObservableList}. De esta manera, se consiguieron solucionar los problemas de concurrencia.

\section{Diseño de la interfaz gráfica}\label{diseño_de_la_interfaz_grafica}

Para el desarrollo de la interfaz gráfica, se ha utilizado la herramienta \emph{SceneBuilder}.

Esta herramienta nos permite crear interfaces gráficas de forma fácil e intuitiva. Una vez realizado el diseño, la herramienta genera un fichero FXML (basado en una estructura XML), el cual será utilizado para mostrar los elementos en pantalla.

En este fichero es importante tener en cuenta los siguientes aspectos:

\begin{itemize}

\item
Debemos de indicarle que controlador se va a encargar de manejar esa vista. Es decir, siempre tiene que existir una relación vista-controlador.
\item
Todos los componentes que vayan a ser utilizados desde el controlador, tiene que tener como mínimo la descripción \emph{fx:id} para su identificación. Esta id tiene que ser única en el documento FXML.
\item
Además, en este fichero se indican las acciones que queremos que tengan los elementos. Por ejemplo, que al pulsar un botón o modificar un texto, el controlador ejecute una función establecida.

\end{itemize}



\section{Testing}\label{testing}

Para realizar las pruebas correspondientes durante el desarrollo de la aplicación y no depender del acceso físico a una máquina o vehículo, se ha desarrollado una pequeña placa la cual simularía una red CAN, en la que se pueden conectar dos nodos, uno para enviar datos al bus y el otro para recibir.

Para ello simplemente se han colocado dos resistencias de 120 \emph{ohmns} (definidas por el estándar del CAN), y unos de tipo DB9, ya que son los más habituales en este tipo de \emph{hardware}.

\imagen{hardwareRedCan}{Placa casera para simular una red CAN}.

\section{Diseño e impresión 3D de una caja protectora}\label{desarrollo_e_impresion_3d_de_una_caja_protectora}

Para la protección de \emph{hardware} desarrollado, se vio la necesidad de crear una caja la cual almacenase dentro la placa, pero que además permitiera la conexión de los elementos necesarios para su funcionamiento.

Para ello se ha utilizado el \emph{software} \emph{SolidWorks} para el diseño de la misma.

FOTO CAJA EN 3D EN EL PC

Una vez concluido el diseño, el fichero 3D ha sido introducido en \emph{Ultimaker CURA}, un \emph{silicer} el cual nos transforma el objeto 3D en un fichero el cual la impresora 3D sea capaz de interpretar.

Para realizar la impresión, se ha utilizado una impresora 3D \emph{BQ Prusa i3 Hephestos} y plástico de tipo ABS (\emph{Acrylonitrile Butadiene Styrene}) el cual resulta más difícil de imprimir, pero es adecuado para un uso intensivo como puede ser el de la industria, teniendo más resistencia a golpes y a temperatura que uno de sus competidores como puede ser el PLA ( Polylactic acid).


FOTO DEL A CAJA TERMINADA


\capitulo{6}{Trabajos relacionados}


Cada fabricante dispone de sus propias herramientas (al igual que muchos de ellos disponen de su propio \emph{hardware}). Estas herramientas, como se ha mencionado en la introducción, tienen un coste muy elevado en la mayoría de los acasos.

Las principales ventajas de este proyecto sobre los otros mencionados anteriormente son:

El hardware desarrollado puede ser producido por cualquier persona, ya que los esquemas y los ficheros necesarios para su producción son open source.

Aplicación multiplataforma, compatible tanto con sistemas Windows como Linux.

No es necesaria la instalación de drivers en el equipo en el que se va a utilizar la herramienta.

Es la única herramienta libre que permite "graficar" los datos que fluyen por el bus CAN, así como su identificación y posterior monitorización.


Las principales desventajas son:

Actualmente solo funciona con un hardware en concreto. Como se menciona en la sección --Lineas futuras de trabajo--, el desarrollo de este proyecto continuaría con la incorporación de un módulo para el soporte de los drivers SocketCAN en sistemas Linux.

A parte, cabe destacar la presencia del proyecto CANalyzat0r, el cual nos permite también el análisis de los datos que pasan por el bus CAN, pero que no nos permite una rápida identificación a través de la visualización de los datos.

Además a través de CRET, se facilita mucho la identificación de las señales, pudiendo etiquetarlas y almacenarlas en el dashboard, y a su vez en una base de datos, para su utilización más adelante. En CANalyzat0r, el etiquetado de los datos identificados tiene que hacerse introduciendo la ID específica, además de que no nos permite separar los bytes del campo de datos.

A continuación se realiza una pequeña comparación entre las herramientas:


\tablaSmall{Comparación entre CRET y otras herramientas del mercado}{l c c c c}{herramientasportipodeuso}
{ \multicolumn{1}{l}{Herramientas} & CRET & CANalyzat0r & PCAN Explorer \\}{
Visualización de los datos en RAW & X & X &X\\
Multiplataforma & X & X &\\
Hardware libre & X & N/A &\\
Representación gráfica de los datos & X & & X\\
Etiquetado con un \emph{click} & X & &\\
Importar y exportar proyectos JSON & X & &\\
} 


\capitulo{7}{Conclusiones y Líneas de trabajo futuras}

\section{Conclusiones}\label{conclusiones}

Una vez concluido el proyecto, es posible extraer las siguientes conclusiones:

Se ha cumplido el objetivo general del proyecto. Ahora es mucho más fácil realizar un análisis de las señales que fluyen por el bus, así como su monitorización.

Para conseguir el desarrollo del proyecto, ha sido necesaria la adquisición de nuevos conocimientos tanto en el ámbito del software, con el uso de JavaFX y la profundización en el lenguaje Java, así como a nivel de \emph{hardware}, en el que se han adquirido conocimientos básicos pero suficientes para alcanzar el objetivo.



\section{Líneas de trabajo futuras}\label{lineas_de_trabajo_futuras}

Es interesante proseguir con el desarrollo de la herramienta y adaptarla para su utilización con los módulos de \emph{Linux} \emph{SocketCAN}.
De esta manera, se conseguiría deshacer de la dependencia que existe actualmente con el \emph{hardware}, para la utilización de la aplicación. \emph{SocketCAN} haría de intermediario entre el \emph{hardware} y el \emph{software} a modo de \emph{driver} de todos aquellos dispositivos que sean soportados.

Otra de las ventajas de realizar esa modificación, sería la posibilidad de crear y utilizar interfaces virtuales dentro del equipo, de manera que ni siquiera sería necesario un \emph{hardware} físico para realizar pruebas.



\bibliographystyle{plain}
\bibliography{bibliografia}

\end{document}
